% Viva / examiner question bank — LiDAR Point Cloud Denoising
% Capstone Project CPG No. 5, Thapar Institute of Engineering and Technology
% Build: tectonic viva_qna.tex
\documentclass[11pt,a4paper]{article}
\usepackage[a4paper,margin=0.9in]{geometry}
\usepackage{amsmath,amssymb}
\usepackage{booktabs}
\usepackage{setspace}
\usepackage[table]{xcolor}
\usepackage{enumitem}
\usepackage[colorlinks=true,linkcolor=black,urlcolor=blue]{hyperref}

\setstretch{1.1}
\setlength{\parindent}{0pt}
\setlength{\parskip}{5pt}

\definecolor{qcol}{HTML}{2A78D6}
\definecolor{warncol}{HTML}{D03B3B}
\definecolor{inkcol}{HTML}{52514E}

% Q / A macros
\newcommand{\Q}[1]{\medskip\noindent\textcolor{qcol}{\textbf{Q.}}\ \textbf{#1}\par\nopagebreak}
\newcommand{\A}[1]{\noindent\hangindent=1.2em\hangafter=0 \textcolor{inkcol}{\textbf{A.}}\ #1\par}
\newcommand{\danger}[1]{\medskip\noindent\textcolor{warncol}{\rule{2.5pt}{9pt}}\ \textit{\textcolor{warncol}{#1}}\par}

\title{\vspace{-1.6cm}\Large\bfseries Examiner Question Bank\\[4pt]
\large LiDAR Point Cloud Denoising for Autonomous Driving Applications}
\author{Capstone Project CPG No.~5 --- Thapar Institute of Engineering and Technology}
\date{August 2026}

\begin{document}
\maketitle
\vspace{-0.5cm}

\begin{center}
\fbox{\begin{minipage}{0.94\textwidth}
\small
\textbf{How to use this document.} Answers are written to be \emph{spoken},
not read out. Where a question has a weak spot, it is flagged in
\textcolor{warncol}{red} --- volunteer those before you are pushed onto them.
Section~\ref{sec:numbers} is a one-page number sheet to revise from.
\end{minipage}}
\end{center}

% ═══════════════════════════════════════════════════════════════════════════
\section{Problem and motivation}

\Q{Why does LiDAR need denoising at all? It is an active sensor.}
\A{Precisely because it is active. It emits 905~nm pulses and times the
returns, so anything airborne that reflects --- rain, fog, snow --- produces
a return. Two distinct failure modes follow. \textbf{Backscatter}: a droplet
reflects the beam before it reaches the target, creating a phantom point in
mid-air. \textbf{Attenuation}: the return is too weak to detect, so a real
point vanishes. Our model reproduces both, and the second turns out to
govern how every result must be read.}

\Q{Why not just use the camera instead?}
\A{Cameras degrade in the same weather and give no direct depth. The point of
this project is that LiDAR degradation is physically structured and therefore
partly predictable --- and so partly reversible --- in a way that generic
image noise is not.}

\Q{What is the actual contribution?}
\A{Three things. A six-method comparison on a common, statistically
controlled protocol. Evidence that reconstruction metrics do not predict
downstream detection quality. And a range-image denoiser that meets the
10~Hz real-time target --- the first learned method in the project to do so,
40$\times$ faster than the graph network.}

% ═══════════════════════════════════════════════════════════════════════════
\section{Data and supervision}

\Q{Are you using real clean/noisy pairs, or adding noise to clean data?}
\A{We add noise --- and that \emph{is} how we get the pairs. Each KITTI scan
is used twice: untouched as the target, corrupted as the input. Because we
generate the corruption we also get per-point ground-truth noise labels,
which is what supervises the DGCNN's classifier and the removal-F1 metric.
You cannot get those labels from real data.}

\Q{Isn't that circular? You are testing whether you can undo your own corruption.}
\A{It is a real limitation and we state it. Three things defend it. First, no
paired real dataset exists --- you cannot record the same scene
simultaneously in rain and not-rain, which is why LISA, Weather-KITTI and
Robo3D all exist. Second, our corruption is physically derived: Mie
extinction $\alpha = 0.003R^{0.6}$ at 905~nm, from published measurements,
not tuned by us. Third, and most important, our headline evaluation is
\emph{downstream detection}, which does not depend on our labels being
correct --- a KITTI-pretrained PointPillars never saw our noise model.}

\danger{Volunteer this before being pushed onto it. Then name the fix: STF,
CADC or WADS --- real measured adverse weather --- is the first thing you
would do next.}

\Q{Is KITTI raw actually ``clean''?}
\A{No, and we should be precise: it is \emph{pseudo}-ground-truth. It is real
sensor data with its own ranging noise and multipath. The models learn to
remove our synthetic corruption layered on top of KITTI's native noise.}

\Q{262 scans from two drives of one session --- is that enough?}
\A{It is enough to train patch-level models --- 28{,}510 training patches ---
but scene diversity is low and both drives are urban daytime. Generalisation
across sensors, cities and lighting is untested and I would not claim it.}

\Q{You do have a self-supervised variant. Does it solve the labelling problem?}
\A{Not yet, and we report that honestly. Pseudo-labels derived from physics
--- range-conditioned intensity anomaly, neighbourhood range incoherence,
edge protection --- score F1 0.484 against the true labels, so roughly half
of what the network is told is wrong. Trained on them the network reaches
F1 0.491: it reproduced its supervision rather than surpassing it. The gap
between 0.49 and the supervised 0.81 is a direct measurement of what
ground-truth labels are worth here.}

% ═══════════════════════════════════════════════════════════════════════════
\section{Noise model}

\Q{Justify your noise model.}
\A{LISA, simplified. An extinction coefficient derived from rain rate gives
two-way transmittance $e^{-2\alpha r}$, and then three channels: backscatter
relocation, lost returns below the detection threshold, and SNR-dependent
range jitter. Gaussian sensor jitter and uniform outliers are added on top.
Severity presets of 10/25/50~mm/h follow the Robo3D three-level convention.}

\Q{Which part of the model matters most?}
\A{The \textbf{deletion} channel. At 25~mm/h it destroys 31\% of returns.
That is why our evaluation has an oracle floor rather than assuming zero
error is reachable.}

\Q{What is the oracle and why does it matter?}
\A{It is what a perfect noise classifier would score --- delete every
injected point, keep every real one. It gets CD \textbf{0.104~m, not zero},
because deleted returns are unrecoverable. Without it you would read
``CD 0.25'' as poor when it has actually closed two-thirds of the
\emph{available} gap. We report an \emph{oracle gap} column for exactly this
reason.}

\Q{Why does intensity matter, and why did you not use it initially?}
\A{Weather backscatter returns far less energy than a solid surface at the
same range --- median 0.051 against 0.150 in our data. The original pipeline
discarded KITTI's intensity channel entirely and worked on geometry alone;
restoring it is what makes IDSOR and the range-image classifier possible.}

% ═══════════════════════════════════════════════════════════════════════════
\section{Methods}

\Q{Why these six methods?}
\A{They span the design space. SOR+ROR is the classical baseline. IDSOR adds
intensity and range adaptivity. Temporal+SOR uses the time axis. PointNet-AE
is a global-bottleneck reconstruction model. DGCNN is a local graph network.
RangeRestoreNet is a 2-D convolutional model on the sensor's native topology.}

\Q{Why does the DGCNN win on reconstruction?}
\A{Local neighbourhood context. EdgeConv lets it distinguish an isolated
backscatter cluster from sparse genuine structure at range, which a global
distance threshold cannot. Its removal precision is 0.969.}

\Q{Why did the autoencoder fail so badly?}
\A{Architecturally it compresses a 256-point patch to a single latent vector
and decodes it with a fully-connected layer. That discards exactly the
high-frequency scan-line structure the detector relies on. It also has no
removal head, so its removal F1 is exactly 0 and its retention exactly 100\%
--- it is a pure displacement model that cannot delete a single spurious
point.}

\Q{What is RangeRestoreNet and why is it so much faster?}
\A{A U-Net over a $64\times2048$ range image with two heads. The speedup is
entirely representational: a rotating LiDAR \emph{is} a 2-D array of (laser
channel, azimuth) samples, so projecting recovers that topology and turns
neighbourhood queries into array indexing. Ordinary 2-D convolutions then
supply the context EdgeConv had to compute by building a $k$-NN graph per
patch. 45~ms versus 1{,}806~ms.}

\Q{What is the cost of that projection?}
\A{Quantisation. At $64\times2048$, 83\% of points win a pixel and 92.8\%
recover their own label after a round trip --- that 92.8\% is the ceiling on
the method's removal accuracy and we report it rather than assume it. The
cause is voxel downsampling replacing the 64 laser rings with grid centroids.
We kept the shared preprocessed cloud anyway, so that all six methods are
compared on identical input.}

\Q{Why is RangeRestoreNet's Chamfer Distance worse than the DGCNN's?}
\A{It is recall-heavy. It finds more of the noise --- removal F1 0.814, the
best of any method --- but removes real points along with it, retaining only
65.6\% of the cloud, below even the oracle's 68.9\%. It is over-aggressive,
and CD punishes that.}

% ═══════════════════════════════════════════════════════════════════════════
\section{Evaluation and statistics}

\Q{How do you know a difference is real and not noise?}
\A{Pairing. Every method is run on the identical corrupted cloud for each
(frame, seed), so differences are paired and a signed-rank test applies.
Geometry uses 30 paired observations, detection \textbf{84}. Every ranking
claim in the report carries a Wilcoxon $p$-value.}

\Q{Why is the detection sample larger than the geometry sample?}
\A{Because we learned it had to be. Detection recall has a standard deviation
of about 0.3 on a mean of 0.4 --- an order of magnitude more variance than
the geometric metrics. At $n=10$ a gap of 0.05 between means is
uninformative. Statistical power has to be established per metric, not once
per project.}

\Q{Isn't one metric enough?}
\A{No, and demonstrating that is part of our contribution. We report Chamfer
Distance, HD$_{95}$, removal precision/recall/F1, retention and downstream
detection, because they disagree. A single-number evaluation would have
called the autoencoder a modest success.}

\Q{Why HD$_{95}$ rather than Hausdorff?}
\A{Strict Hausdorff is the maximum over all residual points, so one survivor
near the 60~m range gate pins it near 30~m for every method. The column
ranked nothing. HD$_{95}$ carries the same information robustly.}

\Q{How do you validate the code itself, as opposed to the models?}
\A{35 unit tests covering behaviour, not just imports: that heavier rain
loses more returns, that ICP recovers a known translation, that HD$_{95}$
ignores the tail strict HD is dominated by, that an emptied cloud returns NaN
rather than a silent infinity, and that the pseudo-labeller does not read the
ground-truth array. Several exist because a defect was found there.}

% ═══════════════════════════════════════════════════════════════════════════
\section{Results}

\Q{Summarise your headline numbers.}
\A{See Section~\ref{sec:numbers}. In one sentence: the DGCNN closes 98.4\% of
the recoverable reconstruction error but runs at 0.55~Hz, while
RangeRestoreNet meets the 10~Hz target at 22~Hz with the best noise-removal
accuracy of any method.}

\Q{Defend the claim that Chamfer Distance is a bad metric.}
\A{Two independent demonstrations. First, the autoencoder improves CD by
8.9\% over the corrupted input while detection recall collapses from 0.259 to
0.062 --- a 76\% loss at $p<10^{-4}$, and the only condition worse than no
denoising at all. That is not a point-count artefact: after our recombination
fix it emits exactly 100\% of the input cardinality. Second, the uncorrected
recombination variant scores a \emph{better} CD (0.525 vs 0.569) while
returning 232\% of the input points --- so the metric is improvable by
duplication as well as by blurring. Across all six methods the CD/recall rank
correlation is only $\rho=-0.60$, $p=0.21$.}

\Q{Your DGCNN scores 0.111 against the oracle's 0.104. What does that mean?}
\A{It is within 0.007~m of a perfect noise classifier --- 98.4\% of the
recoverable gap closed. On detection it is the only method statistically
indistinguishable from the oracle ($p=0.15$); every other method is
significantly worse.}

\Q{Does anything beat the oracle?}
\A{No. And I want to be explicit, because an earlier revision of this work
claimed otherwise. That claim came from ten detection frames with no
significance test. Re-run at 84 paired observations the ranking reverses:
the oracle is best and nothing exceeds it.}

\danger{Do not let this be discovered by the examiner. Say it first --- it
demonstrates exactly the scepticism they are testing for.}

% ═══════════════════════════════════════════════════════════════════════════
\section{Hard questions}

\Q{You found bugs in your own work. Doesn't that undermine the results?}
\A{The opposite --- it is why the current results are trustworthy. Seven
defects are documented with the evidence for each, five confirmed by
measurement rather than inspection. The autoencoder fixes were the risky
ones, because they could have destroyed our headline claim. Re-measured, that
claim got \emph{stronger}.}

\Q{Your temporal filter assumed ego-motion was negligible. Was it?}
\A{No. Measuring it by ICP gave a mean of 0.56~m and a maximum of 0.94~m
against a 0.5~m consistency radius, so on half the frames it exceeded the
threshold and the filter was deleting static structure rather than
backscatter. Compensating improved CD from 0.227 to 0.177 and retained five
percentage points more of the cloud.}

\Q{Both your learned models have a displacement head. Do they work?}
\A{No, and this is one of our two reported negative results. The DGCNN's head
changes mean position error from 0.0804~m to 0.0783~m --- 2.6\% --- and
deleting it entirely leaves Chamfer Distance unchanged. Ours converges to
predicting exactly zero. That is correct behaviour, not undertraining:
weather range jitter is zero-mean, so for a single observation the
error-minimising prediction \emph{is} zero. Undoing it needs a surface prior
over neighbouring returns, which no per-point or per-pixel loss supplies.}

\Q{Your models trained for only 10 epochs. Are they converged?}
\A{The DGCNN is not --- its validation loss was still descending at epoch 10,
so its figures are a floor on what the architecture achieves rather than a
ceiling. We keep the best-on-validation checkpoint, not the last; for the
autoencoder that mattered, since epoch 8 beat epoch 10.}

\Q{You did not meet 10~Hz with the DGCNN.}
\A{Correct --- 0.55~Hz. We met it with RangeRestoreNet at 22~Hz and with the
statistical filters. And I would add a caveat against our own favour: part of
the filters' headroom comes from the weather destroying 31\% of returns, so
the same pipeline is slower in clear conditions, when it is least needed.}

\Q{Quantisation did not help. Why did you try it?}
\A{It was the proposal's Step~5. It shrank the autoencoder 3.4$\times$
(3.34$\rightarrow$0.98~MB), which matters for the embedded Raspberry~Pi
scenario, but gave no speedup because PyTorch's dynamic backend accelerates
Linear layers while our bottleneck is the $k$-NN gather. That is a measured
negative result, not a failure to try.}

\Q{Where is the hardware component from your proposal?}
\A{Deferred --- it needs the physical YDLIDAR~X2. The pipeline is
sensor-agnostic and takes any $(N,3)$ or $(N,4)$ array, so a driver replaces
the KITTI reader with no other change. This is the clearest gap against the
original proposal and I would own it rather than deflect.}

\Q{What would you do with six more months?}
\A{Three things, in order. Validate on measured adverse weather --- STF, CADC
or WADS --- with no retraining, to test whether any of this transfers. Improve
the pseudo-labels with multi-echo returns and temporal consistency, since our
self-supervision is currently bounded by rules that are only 50\% precise.
And attack position restoration properly with an explicit surface-consistency
term, since we have shown a bolted-on displacement head does nothing.}

% ═══════════════════════════════════════════════════════════════════════════
\section{Number sheet}
\label{sec:numbers}

\begin{table}[h]
\centering\small
\setlength{\tabcolsep}{4pt}
\begin{tabular}{lccccc}
\toprule
\textbf{Method} & \textbf{CD (m)} & \textbf{Oracle gap} & \textbf{Removal F1} &
\textbf{Det. recall} & \textbf{Latency}\\
\midrule
Noisy input & 0.572 & 100\% & --- & 0.259 & ---\\
\rowcolor{gray!12}\emph{Oracle} & \emph{0.104} & \emph{0\%} & \emph{1.000} & \emph{0.479} & ---\\
DGCNN & \textbf{0.111} & \textbf{1.6\%} & 0.788 & \textbf{0.452} & 1{,}806~ms\\
IDSOR & 0.174 & 15.0\% & 0.665 & 0.368 & \textbf{49~ms}\\
RangeRestoreNet & 0.215 & 23.8\% & \textbf{0.814} & 0.403 & \textbf{45~ms}\\
SOR+ROR & 0.251 & 31.5\% & 0.628 & 0.392 & 65~ms\\
Temporal+SOR & 0.317 & 45.5\% & 0.567 & 0.389 & 278~ms\\
PointNet-AE & 0.521 & 89.1\% & 0.000 & \textbf{0.062} & 1{,}192~ms\\
\bottomrule
\end{tabular}
\end{table}

\textbf{Numbers worth memorising}

\begin{itemize}[leftmargin=1.4em,itemsep=1pt]
\item \textbf{31\%} --- returns destroyed by the weather model; the reason
the oracle is not zero.
\item \textbf{0.104~m} --- oracle Chamfer Distance, the achievable floor.
\item \textbf{98.4\%} --- fraction of recoverable error the DGCNN closes.
\item \textbf{76\%} --- detection recall lost by the autoencoder while
\emph{improving} CD 8.9\%. $p<10^{-4}$.
\item \textbf{232\%} --- retention at which the buggy recombination scored a
\emph{better} CD.
\item \textbf{40$\times$} --- RangeRestoreNet's speedup over the DGCNN.
\item \textbf{84} --- paired detection observations (28 frames $\times$ 3 seeds).
\item \textbf{0.56~m / 0.94~m} --- mean / max measured ego-motion against a
0.5~m radius.
\item \textbf{2{,}147 of 36{,}033} --- points left uncovered by the original
patch sampling (6\%).
\item \textbf{0.484 $\rightarrow$ 0.491} --- self-supervised F1 against its
own pseudo-labels: the network reproduced its supervision.
\item \textbf{35} --- unit tests passing.
\end{itemize}

\vspace{0.4cm}
\noindent\fbox{\begin{minipage}{0.96\textwidth}
\small
\textbf{Three things to volunteer before being asked.}
\begin{enumerate}[leftmargin=1.4em,itemsep=1pt]
\item The noise is synthetic; real-weather validation is the next step and we
name the datasets.
\item An earlier version of this work claimed removal was not the ceiling;
proper significance testing reversed it, and we corrected it.
\item The hardware component of the proposal is deferred.
\end{enumerate}
Each of these converts a weakness into evidence of judgement --- but only if
you get there first.
\end{minipage}}

\end{document}
